Campaigns 1 and 2 \citep{lyra2026campaign1,lyra2026campaign2} established that KV-cache geometry---the spectral structure of key matrices accumulated during transformer inference---encodes reliable signatures of cognitive processing mode. Campaign 1 discovered category-specific geometry, encoding-native signals, and deception expansion on a single architecture. Campaign 2 confirmed these findings across 6 architecture families, 16 models, and a 140$\times$ parameter range, while adding censorship gradient detection and abliteration analysis.

Both campaigns, however, left a critical question unanswered: \emph{what are the practical limits of this approach?} Establishing that coding prompts produce different geometry from creative writing prompts is scientifically interesting but operationally narrow. For KV-cache analysis to matter for AI safety, we need to know whether it can detect the cognitive states that actually threaten users---deception, harmful compliance, sycophancy---and equally importantly, whether it \emph{cannot} detect states that might be assumed detectable, such as confabulation.

Campaign 3 answers both questions through 28 experiments conducted in rapid succession during the Funding the Commons hackathon. The experimental design is deliberately adversarial: every positive finding is subjected to same-prompt controls, length residualization, token-count ANCOVA, and cross-architecture replication. Every null finding is probed with multiple experimental paradigms before being accepted. Two initially promising results---step-0 encoding-time detection and the contrastive epistemic state signal---were debunked by our own controls and are reported transparently.

\subsection{The Active--Passive Dichotomy}

The central finding of Campaign 3 is a clean partition of cognitive states into \emph{detectable} and \emph{undetectable} categories, determined by whether the state involves active additional processing:

\begin{itemize}[leftmargin=*]
    \item \textbf{Active states} involve the model performing operations beyond standard response generation: representing the truth while constructing a lie (deception), recognizing content it must withhold a response to and declining to answer (refusal), or---in abliterated models---responding to harmful prompts without the suppression circuitry that would normally constrain output (absence of suppression). These states produce detectable geometric signatures because the additional processing (or its removal) creates measurable perturbations in cache structure. The common mechanism is \emph{output suppression}: the model allocating fewer resources per token when withholding.

    \item \textbf{Passive states} involve the model lacking information: generating a plausible-sounding but incorrect answer without internal conflict (confabulation), or processing a user's false claim without recognizing it as false (non-deliberative sycophancy). These states are indistinguishable from normal processing because nothing geometrically different is happening.
\end{itemize}

This dichotomy has direct implications for AI safety monitoring. Internal state inspection (of which KV-cache analysis is one approach) can detect \emph{misbehavior}---cases where a model ``knows better'' and acts against that knowledge---but cannot detect \emph{ignorance}---cases where a model is confidently wrong. Confabulation detection fundamentally requires external verification, not internal monitoring.

\subsection{Summary of Contributions}

\begin{enumerate}[leftmargin=*]
    \item \textbf{Four detection capabilities demonstrated, unified by output suppression}: deception (AUROC $\approx 1.0$ within-model, $\approx 0.86$ cross-model), safety refusal (AUROC $\approx 0.91$), and impossibility refusal (AUROC $\approx 0.94$). Impossibility refusal is \emph{more} detectable than safety refusal, and the two refusal types are geometrically indistinguishable (AUROC $\approx 0.65$)---key evidence that the signal is output suppression, not harm-specific detection. Abliterated models answering harmful prompts show normal (non-suppressed) KV-cache geometry (absence-of-suppression detection, AUROC $\approx 0.88$ vs.\ benign). Extended features (9 total) yield AUROC $\approx 0.95$, though the improvement over 4 features is within CV noise at this sample size. All AUROC estimates at $n = 20$ per class carry $\pm 0.15$--$0.20$ error bars \citep{varoquaux2018cross,hanley1982meaning}. A fifth candidate---sycophancy detection---was refuted: the apparent signal (AUROC 0.9375) was a prompt-length confound, and the length-matched effect size ($d = 0.107$) requires $n \approx 1{,}400$ per condition for 80\% power.

    \item \textbf{Confabulation null established}: five independent experiments across three architectures confirm that confabulation has no detectable KV-cache signature, in either encoding or generation phases. An initially promising contrastive signal ($d = 1.081$) was debunked by a 4-condition control experiment as an information volume artifact.

    \item \textbf{Two-regime mechanistic model}: the KV-cache encodes different information in encoding (structural complexity) versus generation (cognitive intent) phases. A truth axis does not exist in encoding space (cross-model consistency $= -0.046$), but a complexity axis is universal (consistency $= 0.982$).

    \item \textbf{Hardware and scale invariance}: KV-cache geometry is a model property, not a hardware property ($r > 0.999$ between consumer and datacenter GPUs), and category geometry is preserved from 0.6B to 70B parameters ($\rho = 0.83$--$0.90$).

    \item \textbf{Cross-model transfer solved}: logistic regression with standard scaling achieves AUROC $\approx 0.86$ for cross-architecture deception detection, up from $\approx 0.67$ with random forests, by exploiting the linear separability of aggregate cache features.

    \item \textbf{Self-adversarial methodology}: every finding is red-teamed against length confounds, token-count artifacts, prompt-engineering confounds, and cross-architecture consistency. Two promising results (step-0 detection and contrastive epistemic state) were debunked by our own controls and are reported as cautionary examples.
\end{enumerate}
